% Unnumbered research note. No new hereditary count or promoted assertion.
\subsection*{What minimal rank and the canonical saturation comparison actually control}

Work over the algebraically closed ground field in the fixed-anchor
geometric setting.  Let the geometric period be $M$, let $\mathcal A$
be the rank-$M$ generating anchor of the already proved refined-root
construction, and write
\[
 B=q_*\operatorname{End}(\mathcal A),\qquad
 E=q_*\operatorname{Hom}(\mathcal A,\mathcal F),\qquad
 s=\operatorname{rank}(\mathcal F)/M.
\]
Thus $E$ is etale-locally $B^s$.  Suppose the actual canonical stable
construction supplies its full ordinary comparison $Q$ and inclusions
\[
       Q\subset E\subset Q(C),\qquad E=(QB)^{**}.
\]
These are the hypotheses established by the earlier canonical-order
and induced-anchor comparisons, not consequences of rank alone.
All statements below concern actual geometric sheaves; no arbitrary
coefficient-base saturation or relative reflexive-hull flatness is
claimed.

\paragraph{The finite-support quotient.}
Put $T=E/(QB)$.  The equality after double dual, together with the
height-one comparison, says that $T$ has finite support.  The divisor
sandwich gives the more precise annihilation
\[
              \mathcal O_Y(-C)E\subset Q\subset QB,
              \qquad \mathcal O_Y(-C)T=0.
\]
Moreover $T$ is generated by at most $s$ sections as a $B$-module
sheaf.  To see this without assuming that $E$ is globally free, choose
a finite closed subscheme $D$ supporting $T$.  Its connected factors
are local Artin schemes with algebraically closed residue field.  The
etale-local $B$-free frame of $E$ becomes a frame over each such factor:
one can either use a section of the corresponding etale cover, or
lift the residue frame and apply Nakayama to a map between ordinary
free modules of equal rank.  The $s$ images of a frame generate $T$
on that factor.  The central idempotents of the finite scheme $D$
combine these choices into $s$ global sections of $T$.  They define
an actual surjection
\[
                         B^s\longrightarrow T.
\]
This does not lift to a map $B^s\to E$ without further work.  The
sections lie in $H^0(T)$, and their obstruction to lifting through
$E\to T$ lies in $H^1(QB)$.  In the actual canonical comparison
one has $H^0(E)=H^0(Q)=Q_0$: restriction embeds the first space in
the punctured Hom module $Q_0$, while the inclusion of the full
$Q$ gives the reverse inclusion.  Hence also $H^0(QB)=Q_0$, and
the boundary map $H^0(T)\to H^1(QB)$ is injective.  In particular,
no nonzero section of this quotient lifts to a global section of
$E$.  A generator argument must use a genuine twist or a further
construction, rather than presume these quotient generators lift.
Likewise the two controls give no bound for the
ordinary length of $T$.  Already for a regular local ring
$S=k[[z,w]]$, the elementary module
\[
 B=\operatorname{Mat}_M(S),\quad E=B^s,\quad
 T=E/(z,w^N)E
\]
is annihilated by $z$, has $s$ $B$-generators, and has length
$M^2sN$.  This local example is only a test of these two pieces of
information; it is not asserted to arise from a canonical $QB$.

\paragraph{The divisorial defect.}
Let $D_Q$ be the determinant divisor of the inclusion $Q\to E$.
Both ordinary bundles have rank $Mr$, where $r=Ms$.  At each
height-one DVR their relative elementary divisors lie between zero
and the coefficient of $C$.  Therefore
\[
                     0\le D_Q\le Mr\,C.
\]
For the inherited anti-nef comparable cycle $Z$, the existing energy
and intersection bounds give
\[
  0\le-D_Q\cdot Z\le Mr(-C\cdot Z)
                     =O_R(tMr)\le O_R(tr^2).
\]
The ordinary full-bundle degree estimate bounds
$c_1(Q)\cdot Z=O_R(tMr)$, and
$c_1(E)\cdot Z=c_1(Q)\cdot Z+D_Q\cdot Z$ consequently has absolute
value $O_R(tMr)$.  The finite-support quotient $T$ does not change
this divisor.  Hence a bound for $D_Q$ cannot bound $\operatorname{length}T$.
For the order itself the same discussion has rank $M^2$ and gives
the corresponding $O_R(tM^2)$ degree budget.  These are signed degree
bounds, not bounds for both cohomology groups.

\paragraph{Minimal global rank versus residual matrix multiplicity.}
At a component of residue order $m_i$, connectedness of the spectral
torsor forces every character to occur with multiplicity
\[
                              k_i=M/m_i.
\]
Rank $M$ therefore forces rank one on the spectral curve only when
$m_i=M$.  At an unramified component $m_i=1$, the spectral bundle
still has rank $M$.  The generating property and these forced
multiplicities impose no further numerical condition on its coarse
splitting or on the image of the canonical order on that
multiplicity space.

The exact relevant distinction is between the following local
algebra actions.  Let $O$ be a DVR with parameter $z$, fraction field
$K$, and let $V=O^k$.
If a full lattice $\Lambda\subset K^k$ is stable under
$\operatorname{Mat}_k(O)$, then $\Lambda=I^{\oplus k}$ for one
fractional ideal $I$ of $O$.  Indeed the diagonal matrix units give
$\Lambda=\bigoplus I_i e_i$, while the off-diagonal units imply
$I_i=I_j$ for every pair.  Thus the only freedom is a common line
twist.  The same conclusion follows for an order whose reduction
acts as the full matrix algebra: after scaling so that
$\Lambda\subset V$ but $\Lambda\not\subset zV$, its reduction is
a nonzero invariant subspace and hence all of $V/zV$; Nakayama
then gives $\Lambda=V$.

In contrast, the order
\[
              P_{\mathrm{sc}}=O\,I_k+z\operatorname{Mat}_k(O)
\]
has scalar residue action, and every inverse image in $V$ of a
subspace of $V/zV$ is a $P_{\mathrm{sc}}$-stable lattice between
$zV$ and $V$.  This is an actual order action, not merely a divisor
sandwich.  Over a residual curve, a scalar action similarly allows
all elementary kernels of bundle quotients; the saved genus-at-least
two cone example makes the resulting cohomology failure global in
period one and rank two.

This local calculation does not construct a canonical order of
global period $M$ having this residue action.  It identifies the
logical gap precisely: the rank-period theorem gives $k_i=M/m_i$
but does not prove that the canonical residual image is
$\operatorname{Mat}_{k_i}$, nor that its stable lattices have a
bounded splitting.  Thus rank $M$ alone does not supply the premise
which excludes the flag mechanism.  A distinguished anchor or a
separate theorem about its actual residual order action remains
necessary for the proposed deduction.  No new global period-$M$
counterexample is asserted here.

\paragraph{The exact Frobenius comparison which is available.}
For every reflexive module $M_0$ on the fixed covering singularity,
put $F=\operatorname{full}(M_0)$ and
$F^+=\operatorname{full}((F_A^*M_0)^{**})$.  The already proved
source-factorization argument gives
\[
                    F_Y^*F\subset F^+
                         \subset F_Y^*F(pC).
\]
Its two actual factorizations use $x^p,y^p$, and its constants are
independent of the rank and the period.  Actual point-blowup
pullback preserves these maps and inclusions.  Applied to the
canonical orders and their comparison Hom bundles, it gives the
fixed $(p+1)C$ Morita-product conductor used in the finite-object
effectivity theorem.

It does not identify Frobenius pullback of one newly chosen
root-stack anchor with a second relative-full or toric anchor.
Common line twists leave its endomorphism order and canonical
stability unchanged but change that direct comparison.  A proposed
relative-full replacement must therefore either retain the actual
global canonical Hom modules and this Morita context, or prove the
extra compatibility of its normalization with those modules.  The
ordinary full-module comparison is not a proof that arbitrary
canonical-stable root anchors have a prescribed Frobenius lattice.

\paragraph{Status and source scope.}
This is an internal algebraic consequence and scope check of the
already audited rank-period, canonical-order, comparable-cycle and
finite Morita--Frobenius statements.  No new external theorem is
invoked.  It does not give a global small $B$-projective presentation,
a uniform order-valued vanishing, an all-ring object atlas, or a
stronger hereditary representation bound.  The companion
nonminimal-anchor obstruction remains unpromoted and retains its
explicit period-one/rank-two limitation.
